El estudio está cerrado. Estas siete líneas distinguen los ensayos completados de los aspectos no probados al cierre, recogidos en «Lo que permanece sin probar» (sección~\ref{sec:conclusiones}):
\begin{enumerate}
\item \textbf{Severidad del incendio y velocidad de concentración.} Se calculó el dNBR de Sentinel-2 con escenas del 23 de diciembre de 2025 y del 6 de febrero de 2026: \SI{82}{\percent} de las cuencas con señal de quema y \SI{10}{\percent} con severidad moderada-alta o alta. Esa estimación se incorporó al balance con coeficientes de escorrentía supuestos por severidad y al modelo de \citet{gartner2014}, cuya mediana de volumen de mezcla pasó a \SI{148000}{m^3} (sección~\ref{sec:res-incendio}). Permanecen no probados los incrementos de escorrentía y el tiempo de concentración de cada subcuenca según severidad; el ascenso al pico del hidrograma sigue siendo supuesto (\SI{10}{min}).
\item \textbf{Relieve de mayor resolución en el campus.} Las microsimulaciones y m36 reprodujeron la entrada por el claro al patio; m36 calcula \SI{0,17}{m}. Permanece no probada la geometría real de taludes, senderos y accesos mediante un levantamiento LiDAR o fotogramétrico local, ausente en este estudio.
\item \textbf{Separar el agua de lluvia del barro.} No probado: un trazador pasivo que marque el fluido de los hidrogramas permitiría evaluar el perímetro solo con ese aporte, sin contar la escorrentía urbana de la lluvia.
\item \textbf{Edificios porosos y carga gruesa.} La retención del barro en viviendas y la resistencia adicional se ensayaron de forma simplificada (v25 a v27; porosidad verificada en z26). Permanece no probada su calibración con información de daños por vivienda; no se simularon piedras, troncos ni vehículos como carga transportada.
\item \textbf{Validación independiente.} No probada para los ajustes posteriores a la prueba reservada. La prueba con equipamientos sí se ejecutó y fue desfavorable (19/49); no se contrastaron los ajustes con radar Sentinel-1, levantamientos oficiales de depósitos ni marcas de nivel.
\item \textbf{Concentración como variable aleatoria.} Se trató la concentración como incierta mediante $C_v=0{,}20+0{,}30X$ con $X\sim\mathrm{Beta}(3,2)$ (moda 0,40 y media 0,38), y se propagó con cinco nodos de Gauss (q1--q5). Se obtuvo un área esperada de \SI{169}{ha} y un IEC entre 71 y 73 (sección~\ref{sec:res-posterior}). Permanecen no probadas las distribuciones alternativas y la estimación de probabilidades de inundación y percentiles extremos. Se mantienen tres precauciones. La distribución, construida a partir de rangos de la literatura y de la evidencia del evento, es una distribución de creencia y no la distribución estadística real del sedimento, por lo que sus resultados quedan condicionados a ella y su contraste con alternativas (uniforme, Beta(2,2)) no fue probado. La cuadratura estima bien los momentos, pero no las probabilidades de inundación ni los percentiles extremos, porque el mojado y secado introducen discontinuidades, que se pueden tratar interpolando entre simulaciones la concentración umbral a la que cada celda supera un espesor dado o con más simulaciones (muestreo por hipercubo latino). Variar la concentración entre simulaciones representa la incertidumbre de un parámetro, que es un nivel de modelación distinto de la evolución física de la concentración durante el recorrido.
\item \textbf{Evolución de la concentración.} La erosión, el depósito y la incorporación de escombros cambian la concentración a lo largo del recorrido. No probados: un sedigrama de entrada (frente más concentrado y cola más diluida) y un trazador de concentración transportado con el flujo permitirían que la fluencia y la viscosidad varíen con ella. Se distingue así el depósito final de la extensión máxima del flujo, que es la única que se informa al cierre.
\end{enumerate}
